\documentclass[12pt]{article}
\usepackage[utf8]{inputenc}
\usepackage[T1]{fontenc}
\usepackage{amsmath, amssymb}
\usepackage{geometry}
\usepackage{enumitem}
\usepackage{setspace}
\geometry{margin=1in}
\setstretch{1.15}
\setlength{\parindent}{0pt}

\begin{document}

\begin{center}
\Large \textbf{ECON 308 – Final Exam\\Detailed Step-by-Step Solutions}
\end{center}

\vspace{1em}

\section*{Multiple Choice Solutions (1–45)}

\begin{enumerate}[label=\textbf{\arabic*.}]

\item \textbf{Answer: C} ($-MPL_F/MPL_C$; is constant).\\
In the Ricardian model with one factor (labor) and unit labor requirements $a_{LC},a_{LF}$, the PPF is linear. With cloth on the horizontal axis and food on the vertical,
\[
F = \frac{L}{a_{LF}} - \frac{a_{LC}}{a_{LF}} C.
\]
Using $MPL_C = 1/a_{LC}$ and $MPL_F = 1/a_{LF}$, the slope is
\[
\frac{dF}{dC} = -\frac{a_{LC}}{a_{LF}} = -\frac{MPL_F}{MPL_C},
\]
and is constant as you move along the PPF.

\item \textbf{Answer: A} (less; those who are harmed can be compensated by those who gain).\\
From a Kaldor–Hicks perspective, the efficiency criterion focuses on total surplus. Trade raises total surplus, so the distributional effects are \emph{in principle} less central, because gainers could compensate losers and still be better off. In practice this rarely happens, but conceptually that’s why economists emphasize overall welfare.

\item \textbf{Answer: C} (an increase in the price of cloth relative to that of food).\\
In the specific-factors model, labor moves to the sector where the value of its marginal product is higher. If $P_C/P_F$ rises, $P_C\cdot MPL_C$ rises relative to $P_F\cdot MPL_F$, so labor flows into cloth. That raises $L_C$ and lowers $L_F$, increasing the quantity of labor used in cloth.

\item \textbf{Answer: B} (mobile factor).\\
In the standard specific-factors setup there are two goods and three factors: labor, capital specific to one sector, and land specific to the other. Labor can move freely between sectors, so it is the \emph{mobile} factor.

\item \textbf{Answer: E} (two goods; three factors).\\
The specific-factors model has 2 goods and 3 factors: labor ($L$) plus one specific factor for each sector (e.g.\ $K$ and $T$). Hence 2 goods, 3 factors.

\item \textbf{Answer: B} (ambiguous; positive; negative).\\
When a country opens to trade and its export price rises:
\begin{itemize}
\item The specific factor in the export sector gains unambiguously (its real return rises).
\item The specific factor in the import-competing sector loses unambiguously.
\item The mobile factor (labor) has an ambiguous outcome: its real wage in terms of one good rises, but in terms of the other falls.
\end{itemize}

\item \textbf{Answer: A} (in the short run).\\
Specificity of factors is a short-run phenomenon: capital, land, skills, etc.\ are tied to sectors in the short run. In the long run, more factors can move or be repurposed.

\item \textbf{Answer: B} (factor endowments).\\
The Heckscher–Ohlin model explains comparative advantage by relative factor abundance. A country exports the good that uses its abundant factor intensively.

\item \textbf{Answer: C} (the gainers from trade could compensate the losers and still retain gains).\\
H–O implies that trade creates aggregate gains. Even though some factors lose (e.g.\ the scarce factor), total gains are large enough that winners could, in principle, compensate losers and still be better off.

\item \textbf{Answer: D} (PPF shifts outward much more in one direction).\\
Biased growth or biased expansion means the PPF shifts out more for one good than for the other, changing relative supply. That is represented by an outward shift that is steeper in one direction.

\item \textbf{Answer: D} (European capitalists should support U.S.–European free trade).\\
The U.S.\ is relatively land-abundant; the EU is relatively capital-abundant. In H–O, owners of the abundant factor gain from trade, so EU capital owners benefit from U.S.–EU free trade.

\item \textbf{Answer: E} (increase in labor supply raises output of labor-intensive good, reduces capital-intensive good).\\
The Rybczynski theorem in the H–O model: at given goods prices, an increase in the endowment of one factor increases output of the good using that factor intensively and reduces output of the other good.

\item \textbf{Answer: E} (current consumption; decrease).\\
In the intertemporal trade model, a higher real interest rate makes current consumption more expensive relative to future consumption. Households substitute away from current towards future consumption, so current consumption falls.

\item \textbf{Answer: A} (not always guarantee positive changes).\\
An improvement in net commodity terms of trade (export price relative to import price) tends to raise real income, but volume changes or other shocks may offset this. So better terms of trade do not \emph{always} guarantee better outcomes.

\item \textbf{Answer: C} (they create a difference between world and domestic prices).\\
A tariff raises the domestic price above the world price; an export subsidy raises the foreign price and can also change the domestic/export price. Both create a wedge between world and internal prices.

\item \textbf{Answer: B} (more; less; increase).\\
If $P_C/P_F$ rises, the country moves along the PPF to a point with more cloth and less food. Relative supply of cloth, $Q_C/Q_F$, increases.

\item \textbf{Answer: A} (give producers an incentive to export).\\
An export subsidy raises the price producers receive for exports relative to domestic sales, so they allocate more output to export markets.

\item \textbf{Answer: B} (the country could produce either bundle).\\
Points on the PPF represent feasible production combinations given resources and technology. The country can produce either A or B; which one it chooses depends on relative prices and preferences.

\item \textbf{Answer: C} (terms of trade of cloth exporters will improve).\\
$P_C/P_F$ rising means cloth becomes more valuable in world markets relative to food. Countries that export cloth now get more food per unit of cloth, so their terms of trade improve.

\item \textbf{Answer: D} (external economies; natural resources; mobile).\\
Within a country, factors move easily across regions, so interregional trade patterns reflect historical accidents and external economies (agglomeration), rather than fixed factor endowments (like natural resources) as in cross-country trade.

\item \textbf{Answer: C} (increasing returns to scale).\\
If output more than doubles when all inputs double, average cost falls with scale. That is increasing returns to scale.

\item \textbf{Answer: E} (perfect competition unlikely).\\
Internal increasing returns are inconsistent with perfect competition in the long run, because average costs fall with output; firms have market power and prices exceed marginal cost.

\item \textbf{Answer: D} (locate in the same geographic region).\\
External economies come from industry-level clustering: shared suppliers, labor pools, knowledge spillovers. That occurs when firms in the same industry locate near each other.

\item \textbf{Answer: D} (reduce; increased; the industry).\\
External economies of scale are at the industry level. When total industry output rises, each firm’s average cost falls because of shared infrastructure, input suppliers, etc.

\item \textbf{Answer: C} (decrease; decrease).\\
If we consider industry output increasing in the long run:
\begin{itemize}
\item With internal economies, each firm’s average cost falls with its own scale.
\item With external economies, each firm’s average cost falls as the \emph{industry} expands.
\end{itemize}
Either way, higher output is associated with lower average cost.

\item \textbf{Answer: A} (production and consumption distortion effects).\\
A tariff induces:
\begin{itemize}
\item Overproduction by high-cost domestic producers (production distortion),
\item Underconsumption by domestic consumers (consumption distortion).
\end{itemize}
These triangles are the deadweight losses.

\item \textbf{Answer: B} (domestic producers).\\
A tariff raises the domestic price. Domestic producers receive a higher price and expand output, gaining producer surplus. Consumers lose, and the government collects tariff revenue; the main intended beneficiary is the protected domestic industry.

\item \textbf{Answer: C} (raises; lowering).\\
An export subsidy raises the domestic export price and tends to lower the world (importing-country) price because of the extra supply. So prices are higher at home and lower abroad.

\item \textbf{Answer: D} (have dropped substantially).\\
U.S. tariffs were once a main source of government revenue but have fallen dramatically over the post-war period due to multilateral liberalization (GATT/WTO, FTAs).

\item \textbf{Answer: D} (negative effects of the borders).\\
The gravity model shows that borders sharply reduce trade even after controlling for distance and GDP. So BC trades more with other Canadian provinces than with equally distant U.S. states.

\item \textbf{Answer: B} (vertical disintegration).\\
Production broken into many stages in different locations is vertical disintegration / international fragmentation: different firms in different countries handle different stages.

\item \textbf{Answer: B} (increased).\\
World exports as a share of world GDP have trended upward since the 1970s; globalization has deepened trade integration.

\item \textbf{Answer: D} (manufactured goods).\\
China’s exports today are dominated by manufactures (electronics, machinery, consumer goods), even though earlier they relied more on primary products.

\item \textbf{Answer: E} (well-controlled borders).\\
Larger economies, cultural ties, trade agreements, and historical ties all increase trade. Tight or “well-controlled” borders, in the sense of more friction, reduce it.

\item \textbf{Answer: A}.\\
In the Ricardian framework, a sector’s competitiveness depends on its productivity relative to other sectors \emph{and} the wage it has to pay. The relevant condition is $P \ge w/(\text{productivity})$, so both relative productivity and relative wages matter.

\item \textbf{Answer: C} (absolute advantage).\\
If a country uses less labor per unit of a good than another country, it has an absolute advantage in that good. (Comparative advantage is about opportunity costs.)

\item \textbf{Answer: C} ($a_{LC}/a_{LW}$).\\
With unit labor requirements $a_{LC}$ for cheese and $a_{LW}$ for wine, producing 1 unit of cheese uses $a_{LC}$ units of labor. The wine forgone is $a_{LC}/a_{LW}$ units of wine, so
\[
OC(\text{cheese in wine}) = \frac{a_{LC}}{a_{LW}}.
\]

\item \textbf{Answer: C}.\\
Under comparative advantage, a country exports goods it produces at lower opportunity cost and \emph{imports} goods it would produce relatively inefficiently. It thus obtains imports indirectly more efficiently than if it produced them itself.

\item \textbf{Answer: A}.\\
In a many-good Ricardian model, the pattern of exports is determined by relative productivity across sectors within a country: a country exports goods in which it has the highest productivity \emph{relative to its own other goods}, adjusted for wages.

\item \textbf{Answer: B}.\\
Low monetary cost for copper in Rhozundia reflects high absolute productivity or favorable resources, but comparative advantage depends on the opportunity cost relative to other goods. Without knowing that, we cannot be sure.

\item \textbf{Answer: B} (comparative advantage).\\
Comparative advantage is defined by lower opportunity cost. A country that can produce a good at lower opportunity cost than another has a comparative advantage in that good.

\item \textbf{Answer: B}.\\
Immigrants are often positively or negatively selected relative to the native population. Their education, skills, age, and sectoral composition typically differ from the averages of natives.

\item \textbf{Answer: E} (slope $=-MPL_X/MPL_Y$, optimum where slope $=-P_X/P_Y$).\\
With $X$ on the horizontal axis and $Y$ on the vertical, if you move a little labor from $Y$ to $X$, you gain $MPL_X$ units of $X$ and lose $MPL_Y$ units of $Y$. So
\[
\frac{dY}{dX} = -\frac{MPL_X}{MPL_Y}.
\]
At the optimum under free trade, the slope of the PPF equals the slope of the world price line:
\[
-\frac{MPL_X}{MPL_Y} = -\frac{P_X}{P_Y}.
\]

\item \textbf{Answer: depends on the unseen figure (likely B: point c $\to$ point b; increase).}\\
Conceptually:
\begin{itemize}
\item Points on higher indifference curves represent higher utility.
\item A point inside the PPF (like $c$) is inefficient; moving to a point on the PPF that lies on a higher indifference curve (like $b$) raises utility.
\end{itemize}
So moving from an interior or lower-utility point to an efficient point on a higher indifference curve increases welfare. The exact letter depends on the specific labeling in the figure.
\newpage
\item \textbf{Answer: method: opportunity cost $= a_{LC}^*/a_{LF}^*$; numerically depends on table.}\\
Foreign’s opportunity cost of cloth in terms of food equals the labor required for cloth divided by the labor required for food:
\[
OC_{Foreign}(\text{cloth in food}) = \frac{a_{LC}^*}{a_{LF}^*}.
\]
Using the numeric $a_{LC}^*,a_{LF}^*$ from the table on your exam gives one of $\{0.5, 2.0, 6.0, 1.5, 3.0\}$. (You plug in the table values and compute the ratio.)

\end{enumerate}

\vspace{1em}


\subsection*{Problem 1 – External Economies of Scale and Industry Location (China and Mexico)}

\emph{Recall:}
\[
AC_{China} = 40 - 0.02Q_{China}, \qquad
AC_{Mexico} = 42 - 0.01Q_{Mexico}, \qquad
Q^{World} = 300.
\]

Initially:
\[
Q_{China} = 250, \quad Q_{Mexico} = 50.
\]

\textbf{(1) Initial average costs and lower-cost producer}

Compute:
\[
AC_{China}^{(1)} = 40 - 0.02 \times 250 
= 40 - 5 = 35,
\]
\[
AC_{Mexico}^{(1)} = 42 - 0.01 \times 50
= 42 - 0.5 = 41.5.
\]

So at the initial allocation, China has the lower average cost:
\[
AC_{China}^{(1)} = 35 < 41.5 = AC_{Mexico}^{(1)}.
\]
China is the more competitive producer and is the natural location for the industry.

\textbf{(2) Mexico expands to $Q_{Mexico} = 200$}

Suppose Mexico’s domestic market or policy boosts its output to 200. If total world demand remains 300, one simple scenario is:
\[
Q_{Mexico}^{(2)} = 200, \quad Q_{China}^{(2)} = 100.
\]

Compute new average costs:
\[
AC_{Mexico}^{(2)} = 42 - 0.01 \times 200
= 42 - 2 = 40,
\]
\[
AC_{China}^{(2)} = 40 - 0.02 \times 100
= 40 - 2 = 38.
\]

China is still cheaper (38 vs 40), but the gap has narrowed from $6.5$ to $2$.

If Mexico eventually produces even more, say $Q_{Mexico} = 250$ and China $=50$, then
\[
AC_{Mexico} = 42 - 0.01\times 250 = 42 - 2.5 = 39.5,
\]
\[
AC_{China} = 40 - 0.02\times 50 = 40 - 1 = 39.
\]
Costs become very close; with further expansion Mexico could potentially undercut China.

\textbf{(3) Likely production hub}

As Mexico’s industry grows, its average cost falls due to external economies. If Mexico can reach a large enough scale (e.g.\ via domestic demand or subsidies), its $AC$ can approach or even drop below China’s. At that point, firms and future investment may shift towards Mexico, making it a new global hub.

So:
\[
\text{At small }Q_{Mexico}: \text{China is hub. At large }Q_{Mexico}: \text{Mexico can become hub.}
\]

\textbf{(4) Path dependence, multiple equilibria, expectations}

\begin{itemize}
\item \textbf{Path dependence}: The current location (China) reflects historical accidents: early investments, earlier demand, etc. These history-based advantages persist because higher output keeps costs low.
\item \textbf{Multiple equilibria}: There may be one equilibrium with all production in China at low $AC_{China}$ and another with all production in Mexico at low $AC_{Mexico}$. Both are self-consistent outcomes with different histories.
\item \textbf{Role of expectations}: If firms expect that Mexico will become the main production center (e.g.\ because of policy commitments or large domestic market), they invest there, raising $Q_{Mexico}$, lowering $AC_{Mexico}$, and validating the expectation. Expectations can thus be self-fulfilling.
\end{itemize}

This is exactly the logic of external economies of scale: who gets there first (or is perceived as the future hub) can lock in the industry.

\vspace{1em}
\newpage

\subsection*{Problem 2 – Quota and Welfare for a Large Country (United States)}

\emph{Recall:}
\[
Q_D = 1{,}200 - 3P,\quad
Q_S = 2P - 100,
\]
World price $P_W = 200$. Foreign export supply:
\[
P^{Foreign} = 150 + 0.25Q^{Foreign}.
\]
The U.S.\ imposes an import quota of 120 units.

We treat the U.S.\ as large, so its import policy affects the foreign export price.

\textbf{(1) Free-trade price and import quantity}

Under free trade, U.S.\ domestic price equals the world/foreign price $P = P^{Foreign}$ and imports equal U.S.\ import demand, $Q_M^D$.

\textbf{Step 1: U.S.\ import demand as function of $P$}

\[
Q_M^D(P) = Q_D(P) - Q_S(P)
= (1{,}200 - 3P) - (2P - 100)
= 1{,}300 - 5P.
\]

\textbf{Step 2: Equate import demand with foreign export supply}

Let $Q$ be the trade volume. Then
\[
Q = Q_M^D(P) = 1{,}300 - 5P,
\]
and foreign export supply gives
\[
P^{Foreign} = 150 + 0.25Q.
\]
Under free trade $P = P^{Foreign}$, so
\[
P = 150 + 0.25(1{,}300 - 5P).
\]

Solve:
\[
P = 150 + 0.25(1{,}300) - 0.25\cdot 5P
= 150 + 325 - 1.25P
= 475 - 1.25P.
\]
Bring $P$ terms together:
\[
P + 1.25P = 475 \quad\Rightarrow\quad 2.25P = 475,
\]
\[
P^{FT} = \frac{475}{2.25} \approx 211.11.
\]

Then free-trade imports:
\[
Q^{FT} = 1{,}300 - 5P^{FT}
= 1{,}300 - 5\times 211.11
= 1{,}300 - 1{,}055.56
\approx 244.44.
\]

(You can round to $P^{FT}\approx 211$ and $Q^{FT}\approx 244$.)

Check domestic demand and supply:
\[
Q_D^{FT} = 1{,}200 - 3P^{FT} \approx 1{,}200 - 633.33 = 566.67,
\]
\[
Q_S^{FT} = 2P^{FT} - 100 \approx 2\times 211.11 - 100 = 322.22,
\]
\[
Q_D^{FT} - Q_S^{FT} \approx 566.67 - 322.22 = 244.45 \approx Q^{FT}.
\]

\textbf{(2) Under quota: domestic and foreign prices, and quota rents}

Under an import quota of 120, the trade volume $Q^Q$ is fixed at 120 (assuming the quota binds). Let domestic price be $P_Q$ and foreign price $P^{Foreign}_Q$.

\textbf{Step 1: Express imports in terms of $P_Q$}

Imports:
\[
Q^Q = Q_D(P_Q) - Q_S(P_Q)
= (1{,}200 - 3P_Q) - (2P_Q - 100)
= 1{,}300 - 5P_Q.
\]

Set this equal to the quota:
\[
1{,}300 - 5P_Q = 120 \quad\Rightarrow\quad 5P_Q = 1{,}180,
\]
\[
P_Q = 236.
\]

\textbf{Step 2: Foreign export price from export supply}

Trade volume is $Q^Q = 120$, so
\[
P^{Foreign}_Q = 150 + 0.25 \times 120
= 150 + 30 = 180.
\]

\textbf{Step 3: Quota rent}

Each imported unit is bought in the foreign market at $P^{Foreign}_Q = 180$ and sold in the U.S.\ at $P_Q = 236$. The wedge is
\[
P_Q - P^{Foreign}_Q = 236 - 180 = 56.
\]

Quota rents:
\[
\text{Rents} = (P_Q - P^{Foreign}_Q)\times Q^Q
= 56 \times 120 = 6{,}720.
\]

If licenses are given to U.S.\ firms, they receive these rents.

\textbf{(3) Welfare decomposition: production distortion, consumption distortion, rents, ToT effect}

We compare the quota equilibrium to free trade.

\textbf{(a) Domestic quantities under the quota}

At $P_Q = 236$:
\[
Q_D^Q = 1{,}200 - 3\times 236
= 1{,}200 - 708 = 492,
\]
\[
Q_S^Q = 2\times 236 - 100
= 472 - 100 = 372.
\]
Imports:
\[
Q_D^Q - Q_S^Q = 492 - 372 = 120,
\]
consistent with the quota.

\textbf{(b) Changes in domestic quantities vs free trade}

Free trade:
\[
P^{FT}\approx 211.11,\quad Q_D^{FT}\approx 566.67,\quad Q_S^{FT}\approx 322.22,\quad M^{FT}\approx 244.44.
\]

Quota:
\[
P_Q = 236,\quad Q_D^Q=492,\quad Q_S^Q=372,\quad M^Q=120.
\]

So:
\[
\Delta Q_D = Q_D^{FT} - Q_D^Q \approx 566.67 - 492 = 74.67,
\]
(consumers buy less under the quota), and
\[
\Delta Q_S = Q_S^Q - Q_S^{FT} \approx 372 - 322.22 = 49.78,
\]
(domestic producers expand under the quota).

\textbf{(c) Production distortion loss}

The production distortion triangle has base $\Delta Q_S$ and height equal to the domestic price wedge between the quota and free trade, approximately $P_Q - P^{FT}$:
\[
P_Q - P^{FT} \approx 236 - 211.11 = 24.89.
\]

Area:
\[
DWL_{\text{prod}} \approx \frac{1}{2} \times 24.89 \times 49.78 \approx 620. 
\]
(You can round depending on how many decimals you keep.)

\textbf{(d) Consumption distortion loss}

The consumption distortion triangle uses base $\Delta Q_D$ and the same height:
\[
DWL_{\text{cons}} \approx \frac{1}{2} \times 24.89 \times 74.67 \approx 930.
\]

\textbf{(e) Terms-of-trade effect}

Foreign export price under free trade:
\[
P^{FT} \approx 211.11.
\]
Under the quota:
\[
P^{Foreign}_Q = 180.
\]

The U.S.\ now pays less per imported unit; the improvement in terms of trade per unit is roughly
\[
\Delta P^* = P^{FT} - P^{Foreign}_Q \approx 211.11 - 180 = 31.11.
\]

Since the U.S.\ imports 120 units, ToT gain is approximately:
\[
\text{ToT gain} \approx 31.11 \times 120 \approx 3{,}733.
\]

\textbf{(f) Quota rents}

As computed:
\[
\text{Quota rents} = 6{,}720.
\]

If licenses are given to U.S.\ firms, these rents are part of U.S.\ national income. They are not a loss; they are a transfer from foreign producers/consumers to U.S.\ license holders.

\textbf{(g) Net national welfare effect}

From the U.S.\ viewpoint (with quota rents accruing to U.S.\ agents), the net effect is:
\[
\Delta W \approx \text{Quota rents} + \text{ToT gain} - DWL_{\text{prod}} - DWL_{\text{cons}}.
\]

Plugging rough numbers:
\[
\Delta W \approx 6{,}720 + 3{,}733 - 620 - 930 \approx 8{,}903.
\]

So in this stylized large-country example, the quota raises U.S.\ national income relative to free trade, thanks to the terms-of-trade gain and captured rents. (Of course, trading partners lose.)

\textbf{(4) Comparison with an equivalent tariff}

A tariff set so that imports fall to 120 would:
\begin{itemize}
\item Raise the domestic price to the same $P_Q$,
\item Reduce imports to 120,
\item Lower the foreign price to roughly the same $P^{Foreign}_Q$,
\item Create a per-unit wedge equal to the tariff, $t = P_Q - P^{Foreign}_Q \approx 56$.
\end{itemize}

Under such a tariff:
\begin{itemize}
\item The government would collect revenue: $t \times 120 \approx 6{,}720$.
\item The same production and consumption distortions occur (same $P_Q$ and quantities).
\item The terms-of-trade effect is the same (same foreign price).
\end{itemize}

If quota licenses are given free to domestic firms, the \emph{national} welfare outcome of this quota is essentially the same as that of an equivalent tariff: tariff revenue vs.\ quota rents are just different channels of transferring the wedge to domestic agents. If quotas were allocated to foreigners, national welfare would be strictly lower than under an equivalent tariff.

Thus, a tariff that produces the same import volume is:
\begin{itemize}
\item At least as efficient as a quota, and
\item Strictly more efficient if quota rents are not captured by domestic agents.
\end{itemize}

\end{document}
